\section{Literature}\label{sec:literature}

\subsection{What the US evidence established}

The US Trade Adjustment Assistance programme is the world's longest-running
worker-side trade instrument and the source of nearly all causal evidence on
what such instruments do. Three findings organise it. First,
training-centric TAA delivers little: the major federal evaluation
\citep{damicoschochet2012} found four-year earnings impacts of the
training-heavy programme close to zero, with participants' longer training
spells offsetting later earnings gains. Second, income support alone does
not restore earnings: in exposed US labour markets, transfer payments rose
with import competition but offset only a small share of earnings losses,
and arrived mostly through disability and retirement rather than
trade-specific programmes \citep{autor2013, autor2014}. Third, and most
consequentially for design, \emph{wage insurance works through behaviour}:
exploiting the age-50 eligibility discontinuity in Reemployment TAA,
\citet{hyman2024} find eligibility raised employment by 8--17 percentage
points in the two years after displacement and cumulative earnings by
roughly \$18,000 over four years, with effects concentrated in shorter
non-employment spells---large enough that the programme plausibly pays for
itself. The design canon anticipated this: \citet{kling2006},
\citet{lalonde2007} and \citet{kletzerrosen2005} proposed wage-loss
insurance precisely because it subsidises re-employment rather than
non-employment; \citet{wandner2016} reviews the US policy history.

The European counterpart is a cautionary tale about evaluability rather
than a template: the EU Globalisation Adjustment Fund disburses
project-based support after mass redundancies, but its own evaluations
cannot construct counterfactuals \citep{claeyssapir2018}. The UK used the
fund only marginally while a member and left it at Brexit without
replacement.

\subsection{The UK evidence base}

The UK has displacement evidence but no displacement policy literature.
On costs: \citet{hijzen2010} estimate income losses of 14--35 per cent
for displaced UK workers, driven mainly by non-employment spells rather
than re-employment wage penalties; \citet{upwardwright2019} confirm large
and persistent losses with the same decomposition. On trade specifically,
workers in Chinese-import-competing industries lost employment years and
earnings \citep{delyonpessoa2021}, and the deindustrialisation literature
documents the benefit-financed inactivity route through which industrial
job loss became long-term incapacity claims
\citep{beatty2007, beatty2017}. On the fiscal side, the companion study
\citep{ahmadi2026} quantifies the statutory cushioning of a
tariff-calibrated displacement shock---roughly
\TaaCompanionDisplacedCushion\ per cent of gross earnings losses, several
points below the cushioning of an equal-sized diffuse wage cut, with the
displacement-side cushion running through Universal Credit claiming
behaviour and gated by capital rules the survey data cannot fully
exercise. That finding is this paper's motivation: the margin on which
trade shocks are least insured is the margin a designed instrument would
target.

The UK record on instruments is richer than commonly remembered, and its
shape defines this paper's gap precisely. The UK has \emph{had} a
trade-adjustment wage-insurance scheme: under the ECSC readaptation
framework, the Iron and Steel Employees Re-adaptation Benefits Scheme
paid weekly supplements bringing a re-employed redundant steelworker's
earnings up to 90 per cent of previous pay, for periods of 78--130 weeks,
alongside earnings-related unemployment payments for older workers; it
operated in successive statutory forms from the 1970s, helped on the
order of a hundred thousand steelworkers, and ended with the ECSC
Treaty's expiry in 2002 \citep{iserbs1988, iserbsrevocation2012}. The UK
also ran an economy-wide Earnings-Related Supplement to unemployment
benefit from 1966 to 1982. The modern proposal literature revives the
earnings-replacement side only: the Resolution Foundation's Economy 2030
programme designed and costed an earnings-related unemployment benefit
at roughly \pounds 0.4 billion a year---described as ``real wage
insurance,'' though paid while unemployed rather than as a re-employment
top-up \citep{resolutionfoundation2023}; IPPR proposed a repayable
``National Salary Insurance'' \citep{ippr2011}; Policy Exchange a
contributory scheme \citep{policyexchange2022}; and the Institute for
Fiscal Studies has costed the government's own proposed Unemployment
Insurance \citep{ifs2025ui}. None of these is trade-targeted, none
includes a re-employment wage top-up of the kind the US causal evidence
supports and ISERBS embodied, and none reports household-level incidence
through the tax--benefit system. The retraining margin has one recent
structural treatment---\citet{conwell2024} show retraining subsidies
compress the distributional effects of trade shocks, on German data---but
no UK costing. The precise gap, then: no one has designed and costed a
modern UK scheme that is both trade-triggered and wage-insurance-style,
and no proposal of any kind has been priced at the household level. The
institutional review in Section~\ref{sec:institutional} explains why the
gap has survived: ad-hoc intervention has substituted for design since
the statutory scheme lapsed.

\subsection{Where this paper sits}

The exercise combines three traditions. From the US literature it takes
the instrument set and the behavioural evidence; from the UK
microsimulation tradition \citep{dolls2012, brewertasseva2021} it takes
the method---statutory rules applied to representative microdata---and
the discipline of separating mechanical costings from behavioural
scenarios; from the companion study it takes the displacement machinery,
the take-up conventions, and the replication standards. Its contribution
is not causal: nothing here estimates what a UK wage-insurance scheme
would do to re-employment. It is the prior step, absent for the UK and
long available to US policymakers: a designed instrument set with
transparent costings, within which the behavioural questions become
concrete enough to argue about.
